% =====================================================================
%  RAPPORT DE STAGE — OCP
%  Projet : LEAKO (Leaks Survey) — Gestion et inspection des fuites de vapeur
%  Compiler avec : pdflatex (ou Overleaf)
%  Déposer le fichier logo.png (logo OCP) dans le même dossier.
% =====================================================================
\documentclass[12pt,a4paper,oneside]{report}

% ------------------------- PACKAGES -------------------------
\usepackage[utf8]{inputenc}
\usepackage[T1]{fontenc}
\usepackage[french]{babel}
\usepackage{graphicx}
\usepackage{geometry}
\usepackage{xcolor}
\usepackage{titlesec}
\usepackage{fancyhdr}
\usepackage{hyperref}
\usepackage{enumitem}
\usepackage{array}
\usepackage{booktabs}
\usepackage{longtable}
\usepackage{float}
\usepackage{caption}
\usepackage{tikz}
\usepackage{listings}
\usepackage{setspace}
\usepackage{helvet}
\usepackage{amsmath}

% ------------------------- MISE EN PAGE -------------------------
\geometry{top=2.5cm, bottom=2.5cm, left=3cm, right=2.5cm}
\renewcommand{\familydefault}{\sfdefault}
\setstretch{1.3}

% ------------------------- COULEURS OCP -------------------------
\definecolor{ocpvert}{RGB}{0, 122, 92}
\definecolor{ocpvertfonce}{RGB}{0, 82, 62}
\definecolor{gris}{RGB}{90, 90, 90}

% ------------------------- EN-TÊTE / PIED DE PAGE -------------------------
\pagestyle{fancy}
\fancyhf{}
\fancyhead[L]{\small\textcolor{ocpvert}{\textbf{LEAKO — Rapport de stage}}}
\fancyhead[R]{\small\textcolor{gris}{OCP}}
\fancyfoot[C]{\thepage}
\renewcommand{\headrulewidth}{0.4pt}
\renewcommand{\headrule}{\hbox to\headwidth{\color{ocpvert}\leaders\hrule height \headrulewidth\hfill}}

% ------------------------- TITRES -------------------------
\titleformat{\chapter}[display]
  {\normalfont\huge\bfseries\color{ocpvertfonce}}
  {\filleft\Large\chaptertitlename\ \thechapter}
  {20pt}
  {\Huge}
\titleformat{\section}
  {\normalfont\Large\bfseries\color{ocpvert}}
  {\thesection}{1em}{}
\titleformat{\subsection}
  {\normalfont\large\bfseries\color{ocpvertfonce}}
  {\thesubsection}{1em}{}

% ------------------------- LIENS -------------------------
\hypersetup{
  colorlinks=true,
  linkcolor=ocpvertfonce,
  urlcolor=ocpvert,
  citecolor=ocpvert,
}

% ------------------------- DOCUMENT -------------------------
\begin{document}

% =====================================================================
%  PAGE DE GARDE
% =====================================================================
\begin{titlepage}
\begin{center}

\vspace*{1cm}

% Logo OCP (déposer logo.png dans le même dossier)
\includegraphics[width=0.45\textwidth]{logo.png}\\[1.5cm]

{\Huge\bfseries\textcolor{ocpvertfonce}{LEAKO}}\\[0.5cm]
{\Large\bfseries Leaks Survey}\\[1cm]

{\Large Gestion et inspection des fuites de vapeur}\\[0.3cm]
{\large Application mobile, dashboard web et analyse par intelligence artificielle}\\[1.5cm]

\rule{0.6\textwidth}{1pt}\\[1cm]

{\large \textbf{Rapport de stage}}\\[0.5cm]
{\large Présenté à l'encadrant de stage}\\[0.3cm]
{\large \textbf{OCP — Office Chérifien des Phosphates}}\\[2cm]

\begin{minipage}{0.8\textwidth}
\begin{center}
{\large \textbf{Réalisé par :} \underline{\hspace{4cm}}}\\[0.5cm]
{\large \textbf{Encadrant OCP :} \underline{\hspace{4cm}}}\\[0.5cm]
{\large \textbf{Période :} \underline{\hspace{4cm}}}
\end{center}
\end{minipage}

\vfill

{\large \textbf{Année universitaire :} \underline{\hspace{4cm}}}

\end{center}
\end{titlepage}

% =====================================================================
%  REMERCIEMENTS
% =====================================================================
\chapter*{Remerciements}
\addcontentsline{toc}{chapter}{Remerciements}

Je tiens à exprimer ma profonde gratitude à mon encadrant de stage au sein du groupe \textbf{OCP} pour son accompagnement, ses conseils et la confiance qu'il m'a accordée tout au long de ce projet.

Je remercie également l'ensemble de l'équipe qui m'a accueilli et intégré, ainsi que toutes les personnes qui ont contribué, de près ou de loin, à la réalisation de ce travail.

\newpage

% =====================================================================
%  TABLE DES MATIÈRES
% =====================================================================
\tableofcontents
\newpage

% =====================================================================
%  INTRODUCTION GÉNÉRALE
% =====================================================================
\chapter*{Introduction générale}
\addcontentsline{toc}{chapter}{Introduction générale}

Dans le cadre de mon stage au sein du groupe \textbf{OCP}, j'ai été amené à concevoir et développer une application de gestion et d'inspection des fuites de vapeur sur les sites industriels. Les fuites de vapeur représentent un enjeu majeur pour l'industrie : elles engendrent des \textbf{pertes financières importantes}, des \textbf{émissions de CO$_2$} et des \textbf{risques de sécurité} pour les opérateurs.

Le projet \textbf{LEAKO} (Leaks Survey) a pour objectif de fournir aux techniciens de terrain un outil moderne permettant de \textbf{déclarer, suivre et analyser} les fuites de vapeur, tout en calculant automatiquement les pertes financières associées.

Ce rapport présente le projet de manière structurée, en abordant successivement :
\begin{itemize}
    \item le \textbf{contexte} et la \textbf{problématique} du stage ;
    \item l'\textbf{analyse fonctionnelle} (les besoins et fonctionnalités) ;
    \item l'\textbf{analyse non fonctionnelle} (les exigences de qualité) ;
    \item l'\textbf{analyse technique} (l'architecture et les technologies) ;
    \item la \textbf{conclusion} et les perspectives.
\end{itemize}

% =====================================================================
%  CHAPITRE 1 : CONTEXTE ET PROBLÉMATIQUE
% =====================================================================
\chapter{Contexte et problématique}

\section{Présentation de l'organisme d'accueil}

L'\textbf{OCP} (Office Chérifien des Phosphates) est un groupe industriel marocain de renommée mondiale, spécialisé dans l'extraction, la transformation et la commercialisation des phosphates et de leurs dérivés. Le groupe exploite plusieurs sites industriels où la vapeur est utilisée comme vecteur énergétique dans les procédés de production.

\section{Contexte du projet}

Sur les sites industriels, les réseaux de vapeur sont soumis à des contraintes importantes (pression, température, usure) qui peuvent provoquer des \textbf{fuites}. Ces fuites, si elles ne sont pas détectées et réparées rapidement, entraînent :
\begin{itemize}
    \item des \textbf{pertes énergétiques} et financières considérables ;
    \item des \textbf{émissions de gaz à effet de serre} (CO$_2$) ;
    \item des \textbf{risques pour la sécurité} des opérateurs ;
    \item une \textbf{dégradation} des équipements environnants.
\end{itemize}

\section{Problématique}

La gestion des fuites de vapeur était auparavant réalisée de manière \textbf{manuelle et dispersée} : les techniciens notaient les fuites sur papier, sans suivi centralisé, sans calcul précis des pertes et sans outil d'analyse. Ce fonctionnement présentait plusieurs limites :
\begin{itemize}
    \item \textbf{absence de centralisation} des données ;
    \item \textbf{aucun calcul automatique} des pertes financières ;
    \item \textbf{suivi difficile} de l'état de réparation des fuites ;
    \item \textbf{manque d'outils d'analyse} et de reporting pour la prise de décision.
\end{itemize}

\section{Objectifs du projet}

Le projet \textbf{LEAKO} vise à répondre à cette problématique en développant un système complet permettant de :
\begin{enumerate}
    \item \textbf{signaler} les fuites de vapeur depuis le terrain (mobile) ;
    \item \textbf{calculer automatiquement} les pertes financières et les émissions de CO$_2$ ;
    \item \textbf{suivre} l'état de réparation des fuites ;
    \item \textbf{analyser} les données via un dashboard et des rapports ;
    \item \textbf{exploiter l'intelligence artificielle} pour détecter et caractériser les fuites à partir de photos et vidéos.
\end{enumerate}

% =====================================================================
%  CHAPITRE 2 : ANALYSE FONCTIONNELLE
% =====================================================================
\chapter{Analyse fonctionnelle}

L'analyse fonctionnelle consiste à identifier et décrire les \textbf{fonctionnalités} que le système doit offrir pour répondre aux besoins des utilisateurs. Le système \textbf{LEAKO} est composé de trois plateformes complémentaires :
\begin{itemize}
    \item une \textbf{application mobile} (Android) destinée aux techniciens de terrain ;
    \item un \textbf{dashboard web} destiné à la supervision et à l'analyse ;
    \item un \textbf{backend} centralisé qui gère les données et la logique métier.
\end{itemize}

\section{Acteurs du système}

\begin{table}[H]
\centering
\begin{tabular}{@{}ll@{}}
\toprule
\textbf{Acteur} & \textbf{Rôle} \\
\midrule
Technicien de terrain & Signale et suit les fuites via le mobile \\
Superviseur / Responsable & Analyse les données via le dashboard web \\
Administrateur & Gère les paramètres et les utilisateurs \\
\bottomrule
\end{tabular}
\caption{Les acteurs du système}
\end{table}

\section{Identification des besoins fonctionnels}

\subsection{Gestion des utilisateurs et authentification}
\begin{itemize}
    \item Création de compte et connexion sécurisée ;
    \item Gestion du profil utilisateur et changement de mot de passe ;
    \item Authentification par jeton (JWT) pour sécuriser les échanges.
\end{itemize}

\subsection{Gestion des projets collaboratifs}
\begin{itemize}
    \item Création de projets regroupant plusieurs utilisateurs ;
    \item Invitation de membres et gestion des statuts d'invitation ;
    \item Travail collaboratif sur des données centralisées.
\end{itemize}

\subsection{Gestion des campagnes d'inspection}
\begin{itemize}
    \item Création et organisation de campagnes d'inspection par zone ;
    \item Suivi de l'avancement des campagnes ;
    \item Clôture des campagnes terminées.
\end{itemize}

\subsection{Signalement et suivi des fuites}
\begin{itemize}
    \item Saisie complète d'une fuite : tag, date, statut, type de vapeur, pression, diamètre ;
    \item Génération automatique d'un numéro de tag unique ;
    \item Suivi de l'état de la fuite selon quatre statuts : \textit{À réparer}, \textit{En cours}, \textit{Réparée}, \textit{Annulée} ;
    \item Recherche, filtrage et tri des fuites.
\end{itemize}

\subsection{Capture de photos et vidéos}
\begin{itemize}
    \item Capture de photos depuis la caméra ou la galerie ;
    \item Enregistrement de vidéos (durée limitée) ;
    \item Annotation et dessin sur les photos ;
    \item Génération automatique de miniatures.
\end{itemize}

\subsection{Géolocalisation}
\begin{itemize}
    \item Capture de la position GPS de la fuite ;
    \item Fonctionnement de la capture GPS même sans connexion internet ;
    \item Ouverture de la position dans Google Maps.
\end{itemize}

\subsection{Calcul des pertes financières}
\begin{itemize}
    \item Calcul automatique du débit de fuite à partir de la pression et du diamètre ;
    \item Estimation du coût annuel des pertes en \textbf{dirhams} (MAD) ;
    \item Estimation des \textbf{émissions de CO$_2$} associées.
\end{itemize}

\subsection{Analyse par intelligence artificielle}
\begin{itemize}
    \item Analyse automatique des photos et vidéos d'une fuite ;
    \item Détection de la présence d'une fuite ;
    \item Estimation du type, de l'intensité et du diamètre de la fuite ;
    \item Indication du niveau de confiance de l'analyse.
\end{itemize}

\subsection{Communication et collaboration}
\begin{itemize}
    \item Chat par fuite (messages texte) ;
    \item Commentaires audio avec transcription ;
    \item Saisie vocale pour faciliter la saisie sur le terrain.
\end{itemize}

\subsection{Rapports et analyses}
\begin{itemize}
    \item Tableau de bord avec indicateurs clés (KPIs) ;
    \item Métriques financières et taux de réparation ;
    \item Diagrammes et graphiques de répartition ;
    \item Génération et export de rapports en \textbf{PDF}.
\end{itemize}

\subsection{Configuration}
\begin{itemize}
    \item Paramétrage des paramètres de calcul (coût de l'énergie, heures de fonctionnement, facteurs d'émission) ;
    \item Personnalisation des calculs selon les besoins d'OCP.
\end{itemize}

\section{Diagramme des cas d'utilisation}

Le diagramme suivant illustre les principaux cas d'utilisation du système pour les différents acteurs.

\begin{figure}[H]
\centering
\begin{tikzpicture}[
    acteur/.style={draw, circle, minimum size=1cm, fill=ocpvert!15},
    cas/.style={draw, ellipse, fill=white, align=center, minimum width=2.6cm, minimum height=1cm},
    systeme/.style={draw, rectangle, rounded corners, minimum width=12cm, minimum height=9cm, fill=ocpvert!5}
]
% Système
\node[systeme] (sys) at (0,0) {};
\node[above] at (sys.north) {\textbf{Système LEAKO}};

% Acteurs
\node[acteur] (tech) at (-7,2) {\includegraphics[width=0.5cm]{logo.png}};
\node[below] at (tech.south) {\small Technicien};
\node[acteur] (sup) at (-7,-2) {\includegraphics[width=0.5cm]{logo.png}};
\node[below] at (sup.south) {\small Superviseur};

% Cas d'utilisation
\node[cas] (c1) at (-2,3) {Signaler une fuite};
\node[cas] (c2) at (2,3) {Capturer photos/vidéos};
\node[cas] (c3) at (5,1.5) {Géolocaliser};
\node[cas] (c4) at (-2,0) {Suivre l'état};
\node[cas] (c5) at (2,0) {Calculer les pertes};
\node[cas] (c6) at (5,-1.5) {Analyser par IA};
\node[cas] (c7) at (-2,-3) {Consulter les rapports};
\node[cas] (c8) at (2,-3) {Gérer les campagnes};

% Liens acteurs -> cas
\draw[->] (tech) -- (c1);
\draw[->] (tech) -- (c2);
\draw[->] (tech) -- (c3);
\draw[->] (tech) -- (c4);
\draw[->] (tech) -- (c5);
\draw[->] (tech) -- (c6);
\draw[->] (sup) -- (c7);
\draw[->] (sup) -- (c8);
\draw[->] (sup) -- (c4);
\draw[->] (sup) -- (c5);
\end{tikzpicture}
\caption{Diagramme des cas d'utilisation principaux}
\end{figure}

% =====================================================================
%  CHAPITRE 3 : ANALYSE NON FONCTIONNELLE
% =====================================================================
\chapter{Analyse non fonctionnelle}

L'analyse non fonctionnelle définit les \textbf{exigences de qualité} que le système doit respecter, indépendamment des fonctionnalités. Ces exigences garantissent la fiabilité, la sécurité et la performance du système.

\section{Exigences de sécurité}

\begin{itemize}
    \item \textbf{Authentification sécurisée} : accès protégé par mot de passe et jeton JWT ;
    \item \textbf{Confidentialité des données} : les données des utilisateurs et des fuites sont protégées ;
    \item \textbf{Stockage sécurisé} : les jetons de session sont stockés de manière sécurisée sur le mobile ;
    \item \textbf{Communication chiffrée} : les échanges entre les plateformes se font via \textbf{HTTPS}.
\end{itemize}

\section{Exigences de performance}

\begin{itemize}
    \item \textbf{Temps de réponse} : l'application doit répondre rapidement aux actions des utilisateurs ;
    \item \textbf{Gestion des médias} : upload et affichage efficaces des photos et vidéos ;
    \item \textbf{Optimisation} : compression des images et génération de miniatures pour réduire la consommation de données.
\end{itemize}

\section{Exigences de disponibilité et de fiabilité}

\begin{itemize}
    \item \textbf{Disponibilité du service} : le backend est déployé avec redémarrage automatique en cas de panne ;
    \item \textbf{Persistance des données} : les données sont stockées de manière durable et sauvegardées ;
    \item \textbf{Tolérance aux pannes} : le système doit continuer à fonctionner en cas de coupure réseau (capture GPS et saisie possibles hors ligne).
\end{itemize}

\section{Exigences d'ergonomie et d'utilisabilité}

\begin{itemize}
    \item \textbf{Interface intuitive} : l'application mobile est conçue pour être utilisée facilement sur le terrain ;
    \item \textbf{Saisie facilitée} : utilisation de la reconnaissance vocale pour réduire la saisie manuelle ;
    \item \textbf{Feedback visuel} : indicateurs de progression lors des uploads et des analyses.
\end{itemize}

\section{Exigences de portabilité et de compatibilité}

\begin{itemize}
    \item \textbf{Multi-plateforme} : le système est accessible depuis mobile (Android) et depuis un navigateur web ;
    \item \textbf{Compatibilité} : l'application mobile fonctionne sur une large gamme d'appareils Android ;
    \item \textbf{Responsive} : le dashboard web s'adapte aux différentes tailles d'écran.
\end{itemize}

\section{Exigences de maintenabilité et d'évolutivité}

\begin{itemize}
    \item \textbf{Architecture modulaire} : le code est organisé en modules clairs et séparés ;
    \item \textbf{Documentation} : le code et les procédures de déploiement sont documentés ;
    \item \textbf{Évolutivité} : l'architecture permet d'ajouter facilement de nouvelles fonctionnalités.
\end{itemize}

\section{Exigences de déploiement}

\begin{itemize}
    \item \textbf{Déploiement automatisé} : intégration continue et déploiement automatique via GitHub Actions ;
    \item \textbf{Conteneurisation} : le système est déployé via Docker pour garantir la portabilité ;
    \item \textbf{Supervision} : le système est accessible via des domaines sécurisés avec certificats SSL.
\end{itemize}

% =====================================================================
%  CHAPITRE 4 : ANALYSE TECHNIQUE
% =====================================================================
\chapter{Analyse technique}

L'analyse technique décrit l'\textbf{architecture} du système, les \textbf{technologies} utilisées et la manière dont les différents composants interagissent.

\section{Architecture globale}

Le système \textbf{LEAKO} repose sur une \textbf{architecture client-serveur} à trois composants :

\begin{figure}[H]
\centering
\begin{tikzpicture}[
    bloc/.style={draw, rectangle, rounded corners, minimum width=3.2cm, minimum height=1.6cm, align=center, fill=ocpvert!10},
    serveur/.style={draw, rectangle, rounded corners, minimum width=3.2cm, minimum height=1.6cm, align=center, fill=ocpvertfonce!15}
]
\node[bloc] (mobile) at (-4,2) {\textbf{Application mobile}\\Flutter (Android)\\Techniciens de terrain};
\node[bloc] (web) at (4,2) {\textbf{Dashboard web}\\React\\Supervision et analyse};
\node[serveur] (backend) at (0,-2) {\textbf{Backend}\\Spring Boot\\API REST + Base de données};

\draw[<->, thick, ocpvert] (mobile) -- node[left, pos=0.5] {HTTPS} (backend);
\draw[<->, thick, ocpvert] (web) -- node[right, pos=0.5] {HTTPS} (backend);
\end{tikzpicture}
\caption{Architecture globale du système}
\end{figure}

Les deux plateformes frontales (mobile et web) consomment la \textbf{même API REST} fournie par le backend. Cette architecture centralisée garantit la \textbf{cohérence des données} et facilite la maintenance.

\section{Technologies utilisées}

Le tableau suivant récapitule les principales technologies employées pour chaque composant.

\begin{table}[H]
\centering
\begin{tabular}{@{}lll@{}}
\toprule
\textbf{Composant} & \textbf{Technologie} & \textbf{Rôle} \\
\midrule
Backend & Spring Boot (Java 17) & API REST et logique métier \\
 & H2 (mode fichier) & Base de données \\
 & JWT & Authentification \\
 & iText 7 & Génération de rapports PDF \\
 & Google Gemini & Analyse IA des médias \\
\addlinespace
Mobile & Flutter (Dart) & Application Android \\
 & geolocator & Capture GPS \\
 & image\_picker & Photos et vidéos \\
 & speech\_to\_text & Reconnaissance vocale \\
\addlinespace
Web & React + TypeScript & Dashboard web \\
 & Vite & Outil de build \\
 & Tailwind CSS & Interface utilisateur \\
 & Leaflet & Cartes GPS \\
\bottomrule
\end{tabular}
\caption{Technologies utilisées}
\end{table}

\section{Le backend (Spring Boot)}

Le backend constitue le \textbf{cœur du système}. Il expose une API REST sécurisée et gère toute la logique métier.

\subsection{Les entités principales}

Le modèle de données est organisé autour de plusieurs entités :
\begin{itemize}
    \item \textbf{Utilisateur} : les comptes des techniciens et superviseurs ;
    \item \textbf{Projet} : les projets collaboratifs regroupant des utilisateurs ;
    \item \textbf{Campagne} : les campagnes d'inspection par zone ;
    \item \textbf{Fuite} : les fuites de vapeur avec leurs caractéristiques (tag, pression, diamètre, position GPS, coût) ;
    \item \textbf{Photo} : les photos et vidéos associées à une fuite ;
    \item \textbf{Message} : les messages de chat par fuite ;
    \item \textbf{Paramètre global} : les paramètres de calcul (coût de l'énergie, facteurs d'émission).
\end{itemize}

\subsection{Les services}

Le backend est organisé selon le \textbf{pattern Interface + Manager}, ce qui sépare la définition des services de leur implémentation. On distingue notamment :
\begin{itemize}
    \item le service d'\textbf{authentification} (gestion des comptes et des jetons) ;
    \item le service de \textbf{gestion des fuites} ;
    \item le service de \textbf{gestion des médias} (upload, miniatures) ;
    \item le service de \textbf{génération de rapports} PDF ;
    \item le service d'\textbf{analyse IA} (via Google Gemini) ;
    \item le service de \textbf{stockage des fichiers}.
\end{itemize}

\section{L'application mobile (Flutter)}

L'application mobile est destinée aux \textbf{techniciens de terrain}. Elle est développée avec \textbf{Flutter}, un framework moderne permettant de créer des applications performantes.

\subsection{Les écrans principaux}

L'application est organisée autour de plusieurs écrans :
\begin{itemize}
    \item \textbf{Connexion} et création de compte ;
    \item \textbf{Tableau de bord} avec les indicateurs clés ;
    \item \textbf{Gestion des campagnes} (création, détail) ;
    \item \textbf{Gestion des fuites} (création, modification, détail) ;
    \item \textbf{Chat par fuite} ;
    \item \textbf{Rapports et analyses} ;
    \item \textbf{Configuration} des paramètres de calcul ;
    \item \textbf{Gestion des projets} collaboratifs ;
    \item \textbf{Profil} utilisateur.
\end{itemize}

\subsection{Les fonctionnalités techniques}

\begin{itemize}
    \item \textbf{Capture GPS} : obtention de la position précise, même sans connexion internet ;
    \item \textbf{Gestion des médias} : capture, compression et annotation des photos, enregistrement de vidéos ;
    \item \textbf{Reconnaissance vocale} : saisie vocale en français pour faciliter le travail sur le terrain ;
    \item \textbf{Calcul du débit} : application de la \textbf{formule de Napier} pour estimer le débit de fuite :
\end{itemize}

\begin{equation}
\text{Débit (kg/h)} = 0{,}262 \times D^2 \times P_{abs}
\end{equation}

où $D$ est le diamètre de la fuite et $P_{abs}$ la pression absolue. Ce débit permet ensuite de calculer le \textbf{coût annuel} des pertes et les \textbf{émissions de CO$_2$}.

\section{Le dashboard web (React)}

Le dashboard web est destiné à la \textbf{supervision} et à l'\textbf{analyse}. Il est développé avec \textbf{React} et \textbf{TypeScript}.

\subsection{Les pages principales}

\begin{itemize}
    \item \textbf{Tableau de bord} : indicateurs clés, fuites récentes, cercles de progression ;
    \item \textbf{Gestion des fuites} : liste, création, détail, modification ;
    \item \textbf{Gestion des campagnes} : liste, création, détail ;
    \item \textbf{Rapports et analyses} : métriques financières, filtres de période, diagrammes, export PDF ;
    \item \textbf{Configuration} : paramètres globaux ;
    \item \textbf{Application mobile} : QR code et téléchargement de l'APK.
\end{itemize}

\subsection{Les fonctionnalités techniques}

\begin{itemize}
    \item \textbf{Authentification} : gestion des jetons avec rafraîchissement automatique ;
    \item \textbf{Visualisation cartographique} : affichage des fuites sur une carte (Leaflet) ;
    \item \textbf{Graphiques} : diagrammes circulaires et indicateurs de progression ;
    \item \textbf{Export PDF} : génération de rapports personnalisés.
\end{itemize}

\section{L'intelligence artificielle}

L'une des fonctionnalités différenciantes du projet est l'\textbf{analyse automatique des fuites} par intelligence artificielle. À partir d'une photo ou d'une vidéo d'une fuite, le système :
\begin{itemize}
    \item \textbf{détecte} la présence d'une fuite de vapeur ;
    \item \textbf{estime} le type et l'intensité de la fuite ;
    \item \textbf{évalue} le diamètre approximatif ;
    \item \textbf{indique} un niveau de confiance de l'analyse.
\end{itemize}

Cette analyse est réalisée via l'API \textbf{Google Gemini}, intégrée au backend. Elle permet d'\textbf{assister le technicien} dans son évaluation et d'améliorer la \textbf{précision} des données collectées.

\section{Le déploiement}

Le système est déployé sur une \textbf{machine virtuelle} (cloud) et conteneurisé avec \textbf{Docker}.

\begin{figure}[H]
\centering
\begin{tikzpicture}[
    bloc/.style={draw, rectangle, rounded corners, minimum width=3.5cm, minimum height=1.4cm, align=center, fill=ocpvert!10}
]
\node[bloc] (git) at (0,2.5) {\textbf{GitHub}\\Code source};
\node[bloc] (ci) at (0,0) {\textbf{GitHub Actions}\\Intégration continue};
\node[bloc] (vm) at (0,-2.5) {\textbf{Machine virtuelle}\\Docker};
\node[bloc] (back) at (-3,-5) {\textbf{Backend}\\Spring Boot};
\node[bloc] (web) at (3,-5) {\textbf{Dashboard web}\\nginx};

\draw[->, thick, ocpvert] (git) -- (ci);
\draw[->, thick, ocpvert] (ci) -- (vm);
\draw[->, thick, ocpvert] (vm) -- (back);
\draw[->, thick, ocpvert] (vm) -- (web);
\end{tikzpicture}
\caption{Chaîne de déploiement}
\end{figure}

Le processus de déploiement est \textbf{automatisé} : à chaque mise à jour du code sur GitHub, une chaîne d'intégration continue construit les applications et les déploie automatiquement sur le serveur. Le backend et le dashboard web sont exécutés dans des \textbf{conteneurs Docker}, garantissant la portabilité et la reproductibilité de l'environnement.

% =====================================================================
%  CHAPITRE 5 : CONCLUSION
% =====================================================================
\chapter{Conclusion et perspectives}

\section{Bilan du projet}

Le projet \textbf{LEAKO} a permis de développer un système complet de gestion et d'inspection des fuites de vapeur pour le groupe \textbf{OCP}. Ce système offre :
\begin{itemize}
    \item une \textbf{application mobile} ergonomique pour les techniciens de terrain ;
    \item un \textbf{dashboard web} puissant pour la supervision et l'analyse ;
    \item un \textbf{backend centralisé} garantissant la cohérence et la sécurité des données ;
    \item une \textbf{analyse par intelligence artificielle} pour assister les techniciens ;
    \item un \textbf{calcul automatique} des pertes financières et des émissions de CO$_2$.
\end{itemize}

Ce stage m'a permis d'acquérir une \textbf{expérience concrète} dans le développement de systèmes informatiques complets, de la conception à la mise en production, en passant par le développement mobile, web et backend.

\section{Apports personnels}

Ce projet m'a permis de :
\begin{itemize}
    \item \textbf{renforcer} mes compétences en développement mobile (Flutter) et web (React) ;
    \item \textbf{découvrir} le développement backend avec Spring Boot ;
    \item \textbf{apprendre} à intégrer l'intelligence artificielle dans une application ;
    \item \textbf{maîtriser} les outils de déploiement (Docker, intégration continue) ;
    \item \textbf{développer} mon autonomie et ma capacité à gérer un projet de bout en bout.
\end{itemize}

\section{Perspectives}

Plusieurs améliorations pourraient être envisagées à l'avenir :
\begin{itemize}
    \item \textbf{extension} de l'analyse IA à d'autres types de fuites (eau, air comprimé) ;
    \item \textbf{amélioration} des modèles de calcul des pertes ;
    \item \textbf{ajout} de notifications automatiques pour le suivi des réparations ;
    \item \textbf{développement} d'une version iOS de l'application mobile ;
    \item \textbf{intégration} avec les systèmes d'information existants d'OCP.
\end{itemize}

% =====================================================================
%  CONCLUSION GÉNÉRALE
% =====================================================================
\chapter*{Conclusion générale}
\addcontentsline{toc}{chapter}{Conclusion générale}

En conclusion, ce stage au sein du groupe \textbf{OCP} m'a permis de concevoir et de développer une application complète de gestion des fuites de vapeur. Le projet \textbf{LEAKO} répond à un besoin réel de l'entreprise en offrant un outil moderne, centralisé et intelligent pour le suivi des fuites, le calcul des pertes et l'aide à la décision.

Au-delà des aspects techniques, ce stage a été une \textbf{expérience humaine et professionnelle} enrichissante. Il m'a permis de comprendre les enjeux industriels du groupe, de travailler en autonomie et de livrer un produit fonctionnel et déployé.

Je remercie une nouvelle fois mon encadrant et toute l'équipe pour leur accueil et leur accompagnement tout au long de cette période.

\end{document}
